\chapter{Conclusión} \label{chp:conclusion}

%%%%%%%%%%%%%%%%%%%%%%%%%%%%%%%%%%%%%%%%%%%%%%%%%%%%%%%%%%%%%%%%%%%%%%%%%%%%%%%%
%%%%%%%%%%%%%%%%%%%%%%%%%%%%%%%%%%%%%%%%%%%%%%%%%%%%%%%%%%%%%%%%%%%%%%%%%%%%%%%%

La Phase I construye un modelo lógico para predecir \texttt{payment_difficulty_next_month}
con un árbol de decisión interpretable. La validación cruzada mostró que el
árbol irrestricto sobreajusta; limitarlo a profundidad seis y 200 observaciones
por hoja, y conservar cuatro predictores, produce un compromiso más útil entre
generalización y legibilidad. La ROC-AUC media es aproximadamente 0.758 y el
modelo identifica mejor la clase minoritaria cuando se usa el umbral 0.22.

Las reglas aprendidas se apoyan principalmente en el comportamiento reciente de
pago, junto con el plan de asistencia y la edad. Las explicaciones globales y
locales permiten rastrear cada predicción hasta una hoja, pero deben entenderse
como asociaciones predictivas. En particular, la asistencia de reembolso no
puede interpretarse como causa de la dificultad.

Para Phase II quedan pendientes la regresión logística y el Explainable
Boosting Machine, la comparación de sus explicaciones sobre las mismas
observaciones y la generación de predicciones para `test.csv`. También será
necesario revisar la selección del umbral conforme a los costes de error que se
definan para la entrega final.

